\documentclass{ueap-apostila}
\usepackage{float,csquotes}
\usepackage[backend=bibtex,style=authoryear,maxcitenames=2]{biblatex}
\addbibresource{referencias.bib}

\setcounter{tocdepth}{1}
\setcounter{secnumdepth}{1}
\setlength{\headheight}{28pt}
\fancyhead[L]{\small\textcolor{UEAPDarkGreen}{Estatística Aplicada | RENAmazon}}
\fancyhead[R]{\small\textcolor{UEAPDarkGreen}{Unidade \thechapter}}
\fancypagestyle{plain}{\fancyhf{}\fancyfoot[L]{\small Robson Borges de Lima}\fancyfoot[R]{\small\thepage}\renewcommand{\headrulewidth}{0pt}}
\lstdefinestyle{roteiro}{style=ueapRstyle,basicstyle=\ttfamily\fontsize{9}{10.5}\selectfont,
  aboveskip=8pt,belowskip=8pt,xleftmargin=5pt,xrightmargin=5pt,
  frame=l,rulecolor=\color{UEAPDarkGreen},framesep=7pt}
\newcommand{\scriptunidade}[1]{\begingroup\singlespacing\lstinputlisting[style=roteiro]{scripts_r/#1}\endgroup}
\newcommand{\figura}[3]{\begin{figure}[H]\centering\includegraphics[width=#2\linewidth]{figuras/#1}\caption{#3}\end{figure}}
\newcommand{\objetivo}[1]{\begin{objetivos}#1\end{objetivos}}
\newcommand{\quest}[2]{\item[]\textbf{#1.} #2}
\newcommand{\dadosimulados}{\textit{Dados simulados para ensino; não representam levantamentos de campo.}}
\hypersetup{pdftitle={Estatística Aplicada à Ecologia, Biodiversidade e Recursos Naturais Amazônicos},pdfauthor={Robson Borges de Lima},linkcolor=UEAPDarkBlue}

% Adaptação local da capa: preserva faixa, paleta e brasão do template.
% Corrige a posição do texto lateral e acomoda o título mais longo.
\renewcommand{\capaueap}{
\begin{titlepage}\thispagestyle{empty}
\begin{tikzpicture}[remember picture,overlay]
\fill[UEAPDarkGreen] (current page.north west) rectangle ([xshift=1.05cm]current page.south west);
\fill[UEAPGold] ([xshift=1.05cm]current page.north west) rectangle ([xshift=1.15cm]current page.south west);
\node[rotate=90,text=white,font=\bfseries\scriptsize] at ([xshift=.52cm]current page.west) {UEAP\quad |\quad MATERIAL DIDÁTICO};
\end{tikzpicture}
\hspace*{.6cm}\begin{minipage}{.94\textwidth}
\begin{minipage}{.72\linewidth}
{\Large\bfseries\color{UEAPDarkGreen}Universidade do Estado do Amapá}\\[8pt]
{\small Colegiado de Engenharia Florestal\\Programa de Pós-graduação em\\Recursos Naturais Amazônicos}
\end{minipage}\hfill\begin{minipage}{.22\linewidth}\includegraphics[width=\linewidth]{Brasao_UEAP.png}\end{minipage}
\vspace{2cm}
{\small\bfseries\color{UEAPGold}APOSTILA DE APOIO DIDÁTICO}\par\vspace{.45cm}
{\fontsize{25}{30}\selectfont\bfseries\color{UEAPDarkBlue}Estatística Aplicada\\à Ecologia, Biodiversidade\\e Recursos Naturais\\Amazônicos\par}
\vspace{.65cm}{\large\color{UEAPDarkGreen}Fundamentos, métodos estatísticos\\e aplicações práticas em R}\par
\vspace{.6cm}{\color{UEAPGold}\rule{.85\linewidth}{1.2pt}}
\vspace{.6cm}\par{\large\bfseries Prof. Dr. Robson Borges de Lima}\par
\vspace{.2cm}{\small\href{mailto:robson.lima@ueap.edu.br}{robson.lima@ueap.edu.br}}
\vspace{1cm}
\begin{tcolorbox}[colback=UEAPLightGray,colframe=UEAPDarkGreen,boxrule=.5pt,arc=1mm]
\textbf{Aprender com perguntas ecológicas.} Sete unidades com fundamentos, exemplos interpretados, exercícios progressivos e scripts comentados. Dados e análises reproduzíveis para uma formação multidisciplinar.
\end{tcolorbox}
\vspace{.7cm}{\small Macapá, Amapá | 2026\\Versão didática 1.0}
\end{minipage}\end{titlepage}}


\title{Estatística Aplicada à Ecologia, Biodiversidade e Recursos Naturais Amazônicos}
\autor{Robson Borges de Lima}
\disciplina{Estatística Aplicada}
\let\cleardoublepage\clearpage
\begin{document}
\hypersetup{pageanchor=false}
\capaueap
\frontmatter
\fancyhead[R]{\small\textcolor{UEAPDarkGreen}{Orientações e sumário}}
\hypersetup{pageanchor=true}
\renewcommand{\contentsname}{Sumário}
\chapter*{Como estudar com esta apostila}
\addcontentsline{toc}{chapter}{Como estudar com esta apostila}
Esta apostila foi preparada para estudantes de Recursos Naturais Amazônicos com diferentes formações. O ponto de partida é uma pergunta sobre o ambiente, seguida da definição da unidade amostral, da escolha de medidas e da avaliação da incerteza. A programação aparece como instrumento para tornar esse raciocínio verificável.

Cada unidade reúne fundamentos, um exemplo interpretado, exercícios e, ao final, um script completo. Leia a explicação, tente responder às questões e execute o roteiro por blocos. Os comentários identificam as atividades e explicam as escolhas. Os mesmos arquivos \texttt{.R} impressos aqui estão na pasta \texttt{scripts\_r}; assim, a versão digital e a versão executável permanecem sincronizadas.

\begin{atencao}
Todos os conjuntos de dados desta edição são \textbf{simulados}. Servem para ensinar procedimentos, não para descrever a biodiversidade real do Amapá ou estimar efeitos ambientais regionais. Resultados numéricos apresentados foram obtidos executando os arquivos que acompanham o projeto.
\end{atencao}

\textbf{Primeiros passos em R.} Abra \texttt{estatistica-aplicada-ecologia.Rproj} no RStudio. Instale \texttt{ggplot2} uma única vez com \texttt{install.packages("ggplot2")}. Abra o script da unidade e execute uma linha ou seleção com Ctrl+Enter. \texttt{<-} guarda um resultado; \texttt{\$} seleciona uma coluna; \texttt{c()} reúne valores; \texttt{\#} inicia um comentário. Letras maiúsculas e minúsculas são diferentes. Nomes de arquivos devem ser preservados.

Os CSV usam vírgula como separador de campos e ponto decimal. Se seus arquivos usarem ponto e vírgula e vírgula decimal, considere \texttt{read.csv2()}. Confira a importação com \texttt{str()} antes de analisar. Não é necessário saber criar funções ou laços para acompanhar as sete unidades.

\textbf{Para adaptar à dissertação.} A troca do arquivo de dados exige revisar a pergunta, a seleção dos locais, a independência, o esforço, as unidades, os pressupostos e os preditores. Reutilizar um comando não significa reutilizar automaticamente sua justificativa científica. Resultados de modelos devem ser acompanhados das decisões de amostragem e das limitações do estudo.

\textbf{Percurso sugerido.} Unidades 1 e 2: organizar e descrever. Unidade 3: estimar diferenças e sua incerteza. Unidades 4 e 5: compreender experimentos e comparações. Unidade 6: modelar associações. Unidade 7: integrar a análise e redigir resultados. O gabarito ao final oferece respostas e critérios de interpretação.
\tableofcontents
\mainmatter
\fancyhead[R]{\small\textcolor{UEAPDarkGreen}{Unidade \thechapter}}
\chapter{Da pergunta ecológica aos dados}
\objetivo{Distinguir população e amostra; reconhecer variáveis, unidades e réplicas; organizar uma base e delimitar o alcance de uma conclusão.}

\section{Uma pergunta antes de um teste}
Como a diversidade de aves varia com a topografia? Igarapés em paisagens conservadas apresentam mais oxigênio dissolvido? A umidade modifica a decomposição? Essas perguntas envolvem respostas diferentes, coletadas em escalas distintas. A estatística ajuda a resumir a variação, estimar associações e avaliar quanto os dados são compatíveis com hipóteses. Ela não substitui o planejamento de campo.

A \textbf{estatística descritiva} organiza o que foi observado: quantos locais, quais valores, quanta dispersão. A \textbf{inferência} utiliza amostras e modelos para discutir quantidades de uma população ou de um processo. Uma média bem calculada em locais escolhidos apenas pela facilidade de acesso pode ser uma boa descrição desses locais e uma estimativa enviesada da região.

\begin{conceito}
\textbf{População} é o conjunto ao qual a pergunta se refere, delimitado no espaço, no tempo e no critério de inclusão. \textbf{Amostra} é a parte observada. Um \textbf{parâmetro}, como a média populacional $\mu$, descreve a população; um \textbf{estimador}, como $\bar{x}$, é uma regra calculada com a amostra. A estimativa é o valor obtido em uma amostra específica.
\end{conceito}

Se a população de interesse são os igarapés de pequenas bacias de uma região na estação seca, não se deve estender automaticamente a conclusão a grandes rios ou à cheia. Definir a população também ajuda a reconhecer quais ambientes ficaram fora do levantamento.

\section{Variáveis, unidades e escalas}
Uma variável é uma característica registrada por unidade de observação. A \textbf{qualitativa nominal} distingue categorias sem ordem, como tipo de ambiente; a \textbf{ordinal} admite uma ordem, como classes previamente definidas de perturbação. A \textbf{quantitativa discreta} é uma contagem, como riqueza observada; a \textbf{contínua} expressa uma medida, como oxigênio em mg/L. Um número usado como código de local não se torna uma variável quantitativa por estar escrito com algarismos.

\begin{table}[H]\centering\small
\caption{Da pergunta à informação necessária.}
\begin{tabular}{p{.27\linewidth}p{.28\linewidth}p{.35\linewidth}}
\toprule Pergunta & Resposta & Unidade de análise\\\midrule
Diversidade e relevo & Shannon por local & Localidade, com esforço padronizado\\
Qualidade da água & Oxigênio em mg/L & Igarapé, após agregar leituras técnicas\\
Umidade e decomposição & Perda de massa em \% & Recipiente com tratamento independente\\
Clima e riqueza & Número de espécies & Local com informação climática apropriada\\\bottomrule
\end{tabular}\end{table}

As unidades acompanham os valores. Precipitação anual em milímetros não é intercambiável com precipitação de um único mês. A heterogeneidade topográfica deste material é o desvio-padrão da elevação, em metros, dentro de uma janela espacial comum a todos os locais. Não é a elevação média nem uma classe de relevo.

\section{Replicação e independência}
A \textbf{unidade amostral} é a unidade selecionada pelo plano de amostragem; a \textbf{unidade experimental} recebe o tratamento de modo independente. Três leituras do oxímetro no mesmo igarapé são repetições técnicas: ajudam a reduzir erro de medição, mas não representam três igarapés. Quatro pontos de escuta dentro de uma localidade também não garantem quatro réplicas independentes de um atributo medido apenas para essa localidade.

\begin{exemplo}
Em 24 igarapés, três leituras técnicas por igarapé geram 72 registros instrumentais. Ao comparar ambientes na escala do igarapé, pode-se usar uma média por igarapé: $n=24$. Usar $n=72$ como se todas as leituras fossem independentes subestima a incerteza e caracteriza pseudorreplicação.
\end{exemplo}

A independência depende do desenho e do processo ecológico. Locais próximos podem compartilhar condições; visitas sucessivas podem ter autocorrelação. Aumentar a distância não garante independência para qualquer organismo ou variável. Dados agrupados ou repetidos podem exigir modelos que estão além desta introdução.

Na \textbf{amostragem observacional}, o pesquisador registra as condições existentes. Na \textbf{experimentação}, manipula tratamentos e os aloca às unidades. A seleção aleatória de locais favorece representatividade; a alocação aleatória de tratamentos favorece a comparação causal. São operações diferentes. Em um DIC de decomposição, os recipientes devem receber a umidade individualmente: uma única câmara por nível de umidade confundiria tratamento e câmara.

\section{Uma base que permita auditoria}
Use uma linha por unidade e uma coluna por variável; nomes consistentes, identificadores únicos e um dicionário com definições. Registre datas, esforço, protocolo e ausências. Zero é um valor observado; \texttt{NA} indica falta de informação. Em geral, riqueza zero em uma visita não demonstra ausência de espécies no local.

A base principal contém 36 localidades simuladas, 12 por ambiente, com quatro horas de amostragem agregadas por local. A diversidade de Shannon usa logaritmo natural:
\begin{equation}H'=-\sum_{s=1}^{S}p_s\ln p_s,\qquad p_s=n_s/\sum_s n_s.\end{equation}
Aqui, $n_s$ é a abundância da espécie $s$, $S$ é a riqueza observada e termos com $p_s=0$ contribuem zero. $H'$ combina riqueza e distribuição das abundâncias; comunidades com a mesma riqueza podem ter diversidades distintas. Na simulação, $H'$ é uma resposta sintética coerente com $H'\leq\ln S$, sem uma matriz de espécies de campo subjacente.

Esforço igual reduz uma fonte de desigualdade, mas não assegura detectabilidade igual. Chuva, observador e comportamento das aves podem alterar os registros. A organização e exploração inicial no R seguem os princípios apresentados em \textcite{rintro}.

\section{Exercícios da unidade}
\begin{enumerate}
\quest{1 | Conceitual}{Classifique ambiente, riqueza, precipitação e código de local. Diferencie parâmetro de estimativa.}
\quest{2 | Numérico}{Quatro pontos em cada uma de 36 localidades geram 144 registros. Se o clima foi medido na escala da localidade e os pontos foram agregados, qual é o tamanho amostral da análise entre localidades?}
\quest{3 | Em R}{Verifique ausências, duplicação de identificadores e a coerência $H'\leq\ln S$. Explique por que essas verificações não demonstram representatividade.}
\quest{4 | Em R e interpretação}{Exporte local, ambiente e riqueza. Descreva a unidade, o esforço e uma limitação de inferência para essa tabela.}
\end{enumerate}
\clearpage
\section{Roteiro completo em R}
Execute a partir da pasta do projeto. O roteiro responde aos quatro exercícios; comentários conceituais também fazem parte da resolução.
\scriptunidade{01_fundamentos.R}


\chapter{Descrever a variação ecológica}
\objetivo{Calcular e interpretar tendência central, dispersão e quantis; escolher representações gráficas e distinguir variabilidade das observações de precisão da média.}

\section{Qual centro representa os dados?}
A média aritmética distribui igualmente a soma dos valores entre as $n$ observações. A média ponderada dá pesos definidos pelo problema, não escolhidos para produzir um resultado conveniente:
\begin{equation}\bar{x}=\frac{1}{n}\sum_{i=1}^{n}x_i,\qquad
\bar{x}_w=\frac{\sum_i w_i x_i}{\sum_i w_i},\quad w_i\geq0.\end{equation}
O peso precisa ter justificativa. Se três volumes iguais de água têm concentração média 2 mg/L e outro volume igual tem 4 mg/L, a média conjunta é $(3\times2+1\times4)/4=2{,}5$ mg/L. A média simples das duas médias, 3 mg/L, trataria os grupos como tendo o mesmo tamanho. Pesos por esforço não corrigem automaticamente diferenças de detectabilidade ou probabilidades de inclusão.

A \textbf{mediana} divide os dados ordenados em duas metades. Se $n$ é ímpar, ocupa a posição $(n+1)/2$; se é par, usa-se a média dos dois valores centrais. A \textbf{moda} é o valor de maior frequência. Pode haver mais de uma moda, ou nenhum valor com frequência superior aos demais. Para dados contínuos, uma classe modal depende dos intervalos escolhidos.

\begin{exemplo}
Considere nutrientes em cinco amostras de água: $1,2,2,3,12$ mg/L. A média é 4 mg/L, a mediana é 2 mg/L e a moda é 2 mg/L. O evento de 12 mg/L eleva a média e pode representar um pulso ambiental relevante. A mediana descreve melhor uma condição central resistente ao extremo; a média responde à magnitude total dos valores. Nenhuma delas torna a outra incorreta.
\end{exemplo}

Antes de excluir o valor 12, investigue calibração, transcrição, chuva antecedente e protocolo. Uma observação rara pode ser justamente o fenômeno que precisa ser explicado. Para distribuições assimétricas, relatar média e mediana com medidas de dispersão costuma informar mais que selecionar apenas um resumo.

\section{Quanto os valores variam?}
A amplitude total é $A=x_{\max}-x_{\min}$. É intuitiva, mas depende somente dos extremos. A variância e o desvio-padrão amostrais utilizam todas as distâncias em relação à média:
\begin{equation}s^2=\frac{\sum_i(x_i-\bar{x})^2}{n-1},\qquad s=\sqrt{s^2}.\end{equation}
O denominador $n-1$ corresponde aos graus de liberdade depois de estimar uma média. Sob amostragem independente de uma população com variância finita, $s^2$ é um estimador não enviesado de sua variância. A variância tem unidade ao quadrado; o desvio-padrão retorna à unidade original.

\begin{exemplo}
Para $1,2,2,3,12$, os desvios da média 4 são $-3,-2,-2,-1,8$. A soma dos quadrados é $9+4+4+1+64=82$. Logo, $s^2=82/4=20{,}5$ (mg/L)$^2$ e $s\approx4{,}528$ mg/L. A amplitude é 11 mg/L. Esses resultados descrevem forte dispersão, influenciada pelo pulso de nutrientes.
\end{exemplo}

\subsection{Desvio-padrão não é erro-padrão}
O \textbf{DP} descreve a variabilidade entre observações. O \textbf{erro-padrão da média} estima a variabilidade de médias que seriam obtidas em repetições da amostragem, sob o modelo adotado:
\begin{equation}\widehat{EP}(\bar{x})=\frac{s}{\sqrt n}.\end{equation}
No exemplo, $EP\approx2{,}025$ mg/L. A fórmula supõe observações independentes e uma estrutura amostral apropriada. Em um desenho com agrupamento, o cálculo ingênuo pode subestimar a incerteza. Aumentar $n$ tende a reduzir o EP, sem necessariamente reduzir a heterogeneidade ecológica medida pelo DP.

O coeficiente de variação é $CV=100s/\bar{x}$. No exemplo, aproximadamente 113,19\%. Ele expressa dispersão relativa em uma escala de razão, com zero interpretável e média positiva afastada de zero. Não deve ser usado mecanicamente em temperatura Celsius, índices sem zero absoluto ou médias próximas de zero. Classificações universais de CV ``baixo'' ou ``alto'' ignoram diferenças de variável, escala, ambiente e protocolo.

\section{Quartis, decis e percentis}
Um quantil de ordem $p$ localiza uma posição da distribuição. Os quartis correspondem a 25\%, 50\% e 75\%; os decis, a múltiplos de 10\%; os percentis, a múltiplos de 1\%. A mediana é o segundo quartil. A amplitude interquartílica é $AIQ=Q_3-Q_1$, resumindo a dispersão da metade central dos dados.

Para o pequeno exemplo, pelo algoritmo \texttt{type=7} do R, $Q_1=2$, $Q_2=2$, $Q_3=3$ e $AIQ=1$ mg/L. O percentil 90 é 8,4 mg/L: a interpolação produz um valor que não precisa existir na amostra. Outros algoritmos podem gerar diferenças, especialmente com $n$ pequeno. Declare o algoritmo quando a reprodutibilidade do quantil importar.

Um evento acima do percentil 90 de precipitação é alto em relação à distribuição de referência; não constitui, por si, um limiar universal de desastre. É necessário definir o período, o local e a resolução temporal. O mesmo raciocínio vale para percentis de turbidez.

\section{O gráfico deve responder a uma pergunta}
O \textbf{histograma} mostra frequências em intervalos; a escolha de classes afeta sua aparência. A \textbf{densidade} é uma curva suavizada cuja área total é um; a altura não é uma probabilidade pontual. Em variáveis não negativas, a suavização pode estender a curva a valores impossíveis, o que exige cuidado na leitura.

O \textbf{boxplot} resume mediana, quartis e dispersão. Na convenção usual do R, as hastes alcançam observações até 1,5 AIQ além dos quartis. Pontos além das hastes não são automaticamente erros. O \textbf{gráfico de dispersão} mostra pares de valores e permite investigar forma, grupos e observações influentes, sem demonstrar causalidade.

\figura{u2_histograma.pdf}{.80}{Turbidez em 24 igarapés simulados. As frequências dependem da largura das classes; não há evidência causal nesta descrição.}
\figura{u2_boxplot.pdf}{.80}{Oxigênio por ambiente. Cada ponto é um igarapé; o boxplot preserva a visualização da dispersão dentro dos grupos.}

\section{Exercícios da unidade}
\begin{enumerate}
\quest{1 | Numérico}{Calcule média, mediana, moda e média ponderada dos exemplos desta unidade. Explique a diferença entre média e mediana diante do valor 12.}
\quest{2 | Numérico e em R}{Calcule amplitude, variância, DP, EP, CV, quartis, AIQ e percentil 90 de $1,2,2,3,12$. Indique as unidades e o algoritmo de quantis.}
\quest{3 | Em R}{Resuma oxigênio por ambiente e calcule os decis de turbidez da base de igarapés. Defina a unidade de cada resumo.}
\quest{4 | Em R e interpretação}{Produza histograma, boxplot, densidade e dispersão. Discuta assimetria, escolha de classes e por que um gráfico não substitui o desenho amostral.}
\end{enumerate}
\clearpage
\section{Roteiro completo em R}
\scriptunidade{02_descritiva.R}


\chapter{Inferência: diferenças e incerteza}
\objetivo{Formular hipóteses, interpretar valor-p e intervalo de confiança; distinguir grupos independentes de pares; relacionar tamanho de efeito, erros e poder.}

\section{Da amostra à pergunta populacional}
A diferença entre duas médias amostrais é quase sempre diferente de zero, mesmo quando não existe diferença populacional. A inferência avalia a magnitude dessa diferença e sua incerteza, condicionada ao desenho e às suposições. Uma estimativa muito precisa ainda pode estar enviesada por seleção dos locais ou confundimento ambiental.

Para oxigênio dissolvido, defina $\Delta=\mu_C-\mu_M$, em que $C$ indica paisagens conservadas e $M$, mosaicos de uso do solo. A hipótese nula bilateral é $H_0:\Delta=0$; a alternativa é $H_1:\Delta\ne0$. Uma hipótese unilateral, como $H_1:\Delta>0$, precisa de justificativa científica prévia. Escolher a direção depois de observar os resultados altera o procedimento e favorece conclusões indevidas.

\begin{conceito}
O valor-p é a probabilidade, sob $H_0$ e as demais suposições do modelo, de uma estatística de teste tão ou mais extrema que a observada, segundo o critério definido pelo teste. Não é a probabilidade de $H_0$ ser verdadeira, de o resultado ocorrer ``por acaso'' ou de a conclusão estar errada \parencite{wasserstein2016}.
\end{conceito}

O nível de significância $\alpha$ é escolhido antes da análise. Um procedimento com $\alpha=0{,}05$ controla, nas condições em que é válido, a frequência de rejeições indevidas de $H_0$ quando ela é verdadeira. A decisão $p<\alpha$ não informa, sozinha, magnitude, importância ecológica ou qualidade do desenho. Quando $p\geq\alpha$, dizemos que os dados não oferecem evidência suficiente para rejeitar $H_0$; não que a igualdade foi demonstrada.

\section{Normal e t de Student}
A distribuição normal é simétrica e descrita por média $\mu$ e variância $\sigma^2$. Muitas distribuições de erro e de estimadores podem ser aproximadas por ela, mas riqueza, abundância ou variáveis muito assimétricas não se tornam normais por conveniência. O teorema central do limite diz respeito à distribuição de somas ou médias sob condições apropriadas, não à transformação automática dos dados brutos.

Ao estimar a variabilidade com uma amostra, a distribuição t de Student tem caudas mais pesadas que a normal. Os graus de liberdade refletem a informação disponível para estimar a variância. Com graus de liberdade grandes, as duas distribuições se aproximam. Não é necessário decorar tabelas: o R calcula probabilidades e quantis, mas a interpretação continua sendo responsabilidade do pesquisador.

\subsection{O teste t de Welch}
Para dois grupos independentes, o teste de Welch utiliza:
\begin{equation}t=\frac{\bar{x}_C-\bar{x}_M-\Delta_0}
{\sqrt{s_C^2/n_C+s_M^2/n_M}},\qquad \Delta_0=0.\end{equation}
O denominador é o erro-padrão da diferença. Os graus de liberdade são aproximados a partir das variâncias e dos tamanhos de amostra. Ao contrário do teste t com variância combinada, Welch não exige variâncias populacionais iguais. Em R, ele é o padrão de \texttt{t.test(x,y)} \parencite{rwelch}.

O procedimento pressupõe independência e condições adequadas para a aproximação da distribuição da estatística. Com grupos pequenos, assimetria intensa e extremos podem afetar a validade. Avalie o desenho, os gráficos e o mecanismo de geração dos dados, em vez de decidir apenas por um teste de normalidade.

\section{Intervalo de confiança e tamanho de efeito}
Um intervalo de confiança de 95\% é obtido por um procedimento que, em repetições sob o modelo, cobre o parâmetro verdadeiro em 95\% das vezes. O parâmetro é fixo nessa interpretação frequentista. Após calcular um intervalo, não atribuímos a ele uma probabilidade frequentista de 95\% de conter um valor fixo.

Para uma diferença de médias, o intervalo tem a forma:
\begin{equation}\widehat\Delta\;\pm\;t_{0{,}975;\nu}\,
\widehat{EP}(\widehat\Delta).\end{equation}
No teste bilateral correspondente, um IC95\% que exclui zero acompanha a rejeição de $H_0$ a 5\%. Essa correspondência depende de usar o mesmo modelo, a mesma hipótese e o mesmo ajuste de multiplicidade.

\begin{exemplo}
Nos 24 igarapés simulados, as médias foram 6,448 mg/L no ambiente conservado e 5,465 mg/L no mosaico. A diferença foi 0,983 mg/L, com IC95\% de 0,578 a 1,389 mg/L; $t=5{,}047$, $\nu=20{,}468$ e $p\approx0{,}000058$. Há evidência de associação entre ambiente e média de oxigênio neste cenário. A diferença em mg/L é o tamanho de efeito na escala original.
\end{exemplo}

Um limiar de relevância ecológica precisa ser fundamentado para o sistema, a estação e os organismos. Se um limiar previamente definido for 1 mg/L, o intervalo do exemplo inclui efeitos abaixo e acima desse valor: a evidência contra diferença zero não resolve a questão da relevância frente a esse limiar. Não usamos aqui 1 mg/L como padrão ambiental normativo.

\section{Pares, erros e poder}
Quando os mesmos locais são medidos em duas campanhas, calcule $d_i=x_{i,\mathrm{depois}}-x_{i,\mathrm{antes}}$. O teste pareado avalia a média das diferenças:
\begin{equation}t=\frac{\bar d}{s_d/\sqrt n},\qquad gl=n-1.\end{equation}
A normalidade relevante é a das diferenças; os pares devem ser independentes entre si. Ignorar o pareamento perde a estrutura dos dados. Sem controle adequado, uma mudança temporal também pode decorrer de sazonalidade, e não de uma intervenção.

O \textbf{erro tipo I} é rejeitar uma nula verdadeira. O \textbf{tipo II} é não rejeitá-la diante de um efeito específico verdadeiro. O poder, $1-\beta$, é a probabilidade de detectar esse efeito pelo procedimento adotado. Depende da magnitude, variabilidade, replicação, delineamento e $\alpha$. Planeje o tamanho amostral para um efeito cientificamente relevante; não calcule ``poder observado'' apenas para reinterpretar o valor-p já obtido.

\figura{u3_qq.pdf}{.95}{Q--Q por grupo: aproximação à linha sugere compatibilidade com normalidade, mas não a comprova. A independência não é avaliada por esse gráfico.}

\section{Exercícios da unidade}
\begin{enumerate}
\quest{1 | Conceitual}{Escreva $H_0$ e $H_1$ bilaterais para oxigênio e explique por que $p=0{,}03$ não significa 3\% de probabilidade de $H_0$ ser verdadeira.}
\quest{2 | Numérico e em R}{Calcule a diferença Conservado menos Mosaico, o teste de Welch e seu IC95\%. Interprete-os na unidade original.}
\quest{3 | Em R}{Examine Q--Q e Shapiro--Wilk por grupo. Explique por que $p>0{,}05$ em Shapiro não prova normalidade e por que não há teste gráfico de independência.}
\quest{4 | Em R e interpretação}{Analise as duas campanhas da base pareada, preservando a correspondência por local. Discuta a limitação de ausência de grupo-controle.}
\quest{5 | Planejamento em R}{Planeje dois grupos para detectar 0,7 mg/L com DP de 0,6, poder de 80\% e $\alpha=0{,}05$. Arredonde $n$ por grupo para cima e explicite as suposições.}
\end{enumerate}
\subsection*{Oficina de leitura da saída}
No resultado de \texttt{t.test()}, leia primeiro \texttt{mean of x} e \texttt{mean of y}: a ordem define a direção da diferença. Em seguida, localize \texttt{95 percent confidence interval}. Os limites descrevem a diferença populacional compatível com o procedimento; não o intervalo que contém 95\% dos igarapés. Só então consulte \texttt{t}, \texttt{df} e \texttt{p-value} para relacionar a estimativa ao teste.

Escreva duas frases: uma sobre a diferença e sua incerteza; outra sobre a limitação do desenho. Compare com o gabarito apenas depois de redigir. No planejamento amostral, registre qual efeito se deseja detectar, a origem do DP esperado e o número de unidades independentes por grupo. Treze leituras no mesmo igarapé não substituem treze igarapés. Se houver perdas previstas, inclua uma margem justificada no planejamento antes de iniciar a coleta.
\clearpage
\section{Roteiro completo em R}
\scriptunidade{03_inferencia.R}


\chapter{ANOVA, teste F e delineamentos}
\objetivo{Relacionar o delineamento ao modelo; decompor a variabilidade; ler a tabela de ANOVA e diagnosticar os resíduos de DIC e DBC.}

\section{Por que analisar variâncias para comparar médias?}
Para comparar três médias, a ANOVA investiga se a variação entre elas é grande em relação à variação residual dentro dos grupos. O nome se refere à decomposição da variabilidade; a hipótese científica é sobre médias. A hipótese nula de um fator é $H_0:\mu_1=\mu_2=\cdots=\mu_a$. A alternativa diz que pelo menos uma média difere, sem afirmar qual delas ou quantos pares diferem.

No experimento didático, 24 recipientes independentes recebem três níveis de umidade, oito por nível. Todos contêm serapilheira preparada pelo mesmo protocolo e permanecem pelo mesmo período. A resposta é a perda de massa seca, $100(m_0-m_f)/m_0$, em porcentagem. O tempo, a composição inicial e o procedimento de secagem precisam ser padronizados. Os resultados são simulados.

\begin{conceito}
No \textbf{Delineamento Inteiramente Casualizado (DIC)}, os tratamentos são alocados aleatoriamente entre unidades consideradas suficientemente homogêneas para a comparação. A replicação estima a variação não explicada; a casualização reduz a associação sistemática entre tratamento e fatores não controlados. Nenhuma delas elimina erros de medição ou dependência criada pelo protocolo.
\end{conceito}

O recipiente é a unidade experimental porque recebe a umidade independentemente. Sacos de serapilheira dentro do mesmo recipiente são subamostras. Uma câmara única para cada nível de umidade não forneceria replicação independente desse tratamento: qualquer diferença entre câmaras ficaria confundida com a umidade.

\section{Modelo e decomposição}
O modelo de um fator pode ser escrito como:
\begin{equation}Y_{ij}=\mu+\tau_i+\varepsilon_{ij},
\quad i=1,\ldots,a,\quad j=1,\ldots,n_i.\end{equation}
$Y_{ij}$ é a resposta da unidade $j$ no tratamento $i$; $\mu$ representa um nível geral; $\tau_i$ é o efeito do tratamento sob uma restrição de identificação; $\varepsilon_{ij}$ é o erro. No desenho balanceado, pode-se adotar $\sum_i\tau_i=0$. Para os testes clássicos, os erros são independentes, de média zero, variância comum $\sigma^2$ e distribuição normal.

Se $N=\sum_i n_i$, $\bar y_i$ é a média do grupo e $\bar y$ a média geral, então:
\begin{align}
SQ_T&=\sum_i\sum_j(y_{ij}-\bar y)^2,\\
SQ_{\mathrm{trat}}&=\sum_i n_i(\bar y_i-\bar y)^2,\\
SQ_E&=\sum_i\sum_j(y_{ij}-\bar y_i)^2,\\
SQ_T&=SQ_{\mathrm{trat}}+SQ_E.
\end{align}
Cada soma de quadrados mede uma fonte de variação. Os graus de liberdade são $N-1$, $a-1$ e $N-a$, respectivamente. Dividindo cada soma pelo respectivo grau de liberdade, obtemos quadrados médios. Sob a nula, o quadrado médio de tratamentos e o residual estimam a mesma variância; uma separação entre médias tende a elevar o primeiro.

\begin{equation}F=\frac{QM_{\mathrm{trat}}}{QM_E}
=\frac{SQ_{\mathrm{trat}}/(a-1)}{SQ_E/(N-a)}.\end{equation}
A distribuição F de Fisher--Snedecor é não negativa, assimétrica e definida por dois graus de liberdade. Valores grandes fornecem evidência contra a igualdade das médias. O valor-p é a área da cauda superior, sob $H_0$, a partir do F observado.

\begin{table}[H]\centering\small
\caption{ANOVA do DIC simulado. Somas e quadrados médios arredondados; perda em porcentagem.}
\begin{tabular}{lrrrrr}\toprule
Fonte & gl & SQ & QM & F & valor-p\\\midrule
Umidade & 2 & 1248,1 & 624,1 & 34,18 & $2{,}49\times10^{-7}$\\
Resíduo & 21 & 383,4 & 18,3 & --- & ---\\
Total & 23 & 1631,6 & --- & --- & ---\\\bottomrule
\end{tabular}\end{table}

Arredondamentos podem fazer as parcelas impressas diferirem ligeiramente do total. No cálculo sem arredondamento, a identidade é exata. O teste indica evidência contra médias iguais, mas não identifica sozinho os pares distintos. A unidade seguinte examina essas comparações.

\section{Diagnosticar é confrontar o modelo com os dados}
Resíduos são diferenças entre observações e valores ajustados. Examine resíduos versus ajustados: padrões de curvatura ou abertura em funil sugerem problemas. O Q--Q investiga compatibilidade aproximada da distribuição residual com a normal. Resíduos por tratamento e o teste de Fligner auxiliam na avaliação das variâncias. Com amostras pequenas, a ausência de rejeição não comprova homogeneidade.

\figura{u4_diagnostico.pdf}{.84}{Diagnóstico do DIC. Os dois painéis superiores avaliam forma e distribuição; os inferiores ajudam a identificar dispersão e influência. Nenhum painel confirma independência.}

Independência vem principalmente da alocação e da condução experimental. Um gráfico por posição ou ordem de medição pode revelar padrões, mas a ausência deles não valida um protocolo inadequado. Investigue observações influentes e documente decisões. Não elimine pontos apenas para obter $p<0{,}05$.

Se as variâncias forem muito distintas, a ANOVA de Welch pode ser pertinente em comparações de um fator com grupos independentes. Ela não corrige pseudorreplicação. Respostas limitadas entre 0 e 100\%, sobretudo próximas das bordas, podem exigir outra estrutura de distribuição. Neste experimento, os valores ficam suficientemente afastados das bordas para explorar o modelo normal, com diagnóstico obrigatório.

\section{Quando o ambiente sugere blocos}
No \textbf{Delineamento em Blocos Casualizados (DBC)}, unidades são agrupadas em blocos de condições relativamente semelhantes; todos os tratamentos aparecem em cada bloco e são casualizados dentro dele. Imagine seis blocos locais distribuídos ao longo de um gradiente topográfico. Cada bloco recebe três unidades, uma para cada tratamento. Bloqueamento reduz variação residual quando explica heterogeneidade relevante \parencite{quinn2002}.

\begin{equation}Y_{ij}=\mu+\tau_i+\beta_j+\varepsilon_{ij}.\end{equation}
O termo $\beta_j$ representa o efeito do bloco. Para $a=3$ tratamentos e $b=6$ blocos, temos 2 gl de tratamento, 5 de bloco e $(a-1)(b-1)=10$ de resíduo. O modelo introdutório é aditivo. Com uma observação por combinação, não se estima separadamente uma interação tratamento--bloco; variação dessa natureza pode aparecer no resíduo.

Não se deve formar blocos depois de olhar a resposta para melhorar a significância. A identificação dos blocos pertence ao planejamento. Os blocos deste exemplo são unidades locais distintas; não se pretende estimar um ``efeito causal de topografia'' por meio de seus rótulos.

\section{Exercícios da unidade}
\begin{enumerate}
\quest{1 | Conceitual e em R}{Identifique tratamento, resposta e unidade experimental no DIC. Confira oito recipientes por nível e ajuste o modelo.}
\quest{2 | Numérico e em R}{Recalcule $SQ_T$, $SQ_E$, $SQ_{\mathrm{trat}}$ e F. Explique os graus de liberdade 2, 21 e 23.}
\quest{3 | Em R e interpretação}{Examine os quatro gráficos de diagnóstico, a posição dos recipientes e Fligner. Indique quais suposições cada informação pode ou não avaliar.}
\quest{4 | Em R}{Ajuste o DBC com bloco e tratamento. Confira o balanceamento, os graus de liberdade e explique por que a interação não é estimada separadamente.}
\end{enumerate}
\clearpage
\section{Roteiro completo em R}
\scriptunidade{04_anova.R}


\chapter{Comparar médias com critério}
\objetivo{Distinguir teste global de comparações específicas; interpretar Tukey, contrastes e intervalos; manter a magnitude ecológica no centro da conclusão.}

\section{Uma pergunta global e perguntas específicas}
O teste F responde se todas as médias são compatíveis com a igualdade sob o modelo. Para três tratamentos, há três comparações par a par: Média--Baixa, Alta--Baixa e Alta--Média. Uma rejeição global não significa que todos os pares diferem. Comparações planejadas expressam perguntas formuladas antes de observar os dados; análises pós-hoc exploram diferenças dentro de uma família de comparações definida e precisam de interpretação compatível.

Fazer vários testes sem ajuste eleva a probabilidade de pelo menos um falso positivo na família. Para $m$ testes independentes de nível $\alpha$, essa probabilidade seria $1-(1-\alpha)^m$ quando todas as nulas fossem verdadeiras. Comparações que compartilham médias não são independentes; a expressão ilustra o problema, mas não calcula seu risco exato nesses casos.

\begin{conceito}
O controle do erro familiar procura limitar a probabilidade de ao menos uma rejeição indevida em uma família de hipóteses. Defina a família pela pergunta científica. A multiplicidade não desaparece ao selecionar apenas as comparações com valores-p menores.
\end{conceito}

Uma estratégia didática é começar pelo F global e então examinar Tukey. Contudo, o procedimento de Tukey já foi construído para controlar a família de comparações e não depende matematicamente de uma rejeição prévia do F. Relate a estratégia escolhida sem transformar regras operacionais em exigências universais.

\section{Teste de Tukey e intervalos simultâneos}
O método de Tukey compara todos os pares de médias usando a distribuição da amplitude studentizada. Para o desenho balanceado, com $n$ unidades por grupo, a diferença mínima associada ao procedimento tem a forma:
\begin{equation}D_{\mathrm{Tukey}}=q_{1-\alpha;a,\nu}\sqrt{QM_E/n},\end{equation}
em que $a$ é o número de grupos e $\nu$ os graus de liberdade residuais. Os intervalos das diferenças são $\bar y_i-\bar y_j\pm D_{\mathrm{Tukey}}$. O uso do mesmo erro residual pressupõe a adequação do modelo com variância comum. O R incorpora uma adaptação para tamanhos amostrais diferentes; desequilíbrios fortes e heterogeneidade exigem atenção adicional \parencite{rtukey}.

\begin{table}[H]\centering\small
\caption{Diferenças de perda de massa no DIC, em pontos percentuais, com IC95\% simultâneos de Tukey.}
\begin{tabular}{lrrr}\toprule
Comparação & Diferença & Limite inferior & Limite superior\\\midrule
Média menos Baixa & 7,775 & 2,390 & 13,160\\
Alta menos Baixa & 17,624 & 12,239 & 23,009\\
Alta menos Média & 9,849 & 4,464 & 15,234\\\bottomrule
\end{tabular}\end{table}

Os três intervalos excluem zero. O contraste mais amplo, Alta menos Baixa, indica maior perda média em umidade alta, de aproximadamente 17,6 pontos percentuais. Trata-se de uma diferença entre porcentagens, não de aumento relativo de 17,6\%. Um aumento relativo exigiria dividir a diferença pela média de referência.

A interpretação ecológica é restrita ao período, substrato e condições representados pelo experimento. Maior perda de massa sugere decomposição mais intensa, mas não mede diretamente toda a ciclagem de nutrientes, a mineralização ou a comunidade decompositora. Um mecanismo plausível precisa ser distinguido de uma variável efetivamente medida.

\section{Médias, intervalos e letras}
Um gráfico de médias com intervalos de confiança informa localização e precisão. Os IC95\% individuais de cada média não são os mesmos que os intervalos simultâneos das diferenças de Tukey. A sobreposição entre IC individuais não determina, por inspeção, o resultado do teste para uma diferença.

\figura{u5_medias.pdf}{.85}{Médias de perda de massa e IC95\% individuais calculados com o quadrado médio residual do DIC. Para as comparações simultâneas, utilize a tabela de diferenças de Tukey.}

Letras de agrupamento resumem decisões binárias, mas ocultam magnitude e incerteza. Compartilhar uma letra indica que o procedimento não encontrou evidência suficiente de diferença naquele nível; não demonstra equivalência. Não compartilhar letras não informa se a diferença é ecologicamente importante. Nesta apostila, priorizamos estimativas e intervalos; as letras não são necessárias para compreender os resultados.

\section{Contrastes: a hipótese em forma de diferença}
Um contraste é uma combinação linear de médias com coeficientes que somam zero:
\begin{equation}L=\sum_i c_i\mu_i,\qquad\sum_i c_i=0.\end{equation}
Para Alta menos Baixa, na ordem Baixa, Média, Alta, use $c=(-1,0,1)$. Sob o modelo independente com variância comum:
\begin{equation}\widehat{EP}(\widehat L)=\sqrt{QM_E\sum_i c_i^2/n_i},
\qquad t=\widehat L/\widehat{EP}(\widehat L).\end{equation}
No DIC, o erro-padrão desse contraste é aproximadamente 2,136 pontos percentuais. Um único contraste previamente definido pode ser analisado como uma pergunta única. Se vários contrastes integram a mesma família, planeje o controle de multiplicidade. Não descreva como planejado um contraste escolhido por apresentar a maior diferença observada.

O pacote \texttt{emmeans} é útil para médias ajustadas e contrastes em modelos mais elaborados. Para evitar dependências adicionais nesta introdução, o script explicita o cálculo do contraste no desenho balanceado. A fórmula não deve ser transferida sem revisão para modelos mistos, dados dependentes ou outros estimadores de variância.

\section{Exercícios da unidade}
\begin{enumerate}
\quest{1 | Conceitual e em R}{Explique a diferença entre rejeitar a nula global e identificar um par distinto. Ajuste novamente o DIC.}
\quest{2 | Numérico e em R}{Obtenha Tukey. Para Alta--Baixa, interprete diferença, intervalo e valor-p ajustado; explique a unidade ``pontos percentuais''.}
\quest{3 | Em R e interpretação}{Produza o gráfico das médias com IC individuais. Explique por que ele não substitui os IC simultâneos das diferenças.}
\quest{4 | Em R}{Calcule o contraste Alta menos Baixa com a fórmula apresentada. Explique em que situação pode ser tratado como uma única hipótese planejada.}
\end{enumerate}
\subsection*{Oficina de interpretação das comparações}
Ao ler cada linha de Tukey, identifique a direção da subtração antes de interpretar o sinal. Depois, verifique se o intervalo inclui zero e avalie sua amplitude. Por fim, discuta a relevância ecológica da diferença. Uma ordem de leitura útil é: \textbf{qual diferença, de quanto, com quanta incerteza e para qual pergunta?}

Para Alta menos Baixa, redija o resultado em pontos percentuais e explique qual processo foi medido. Compare essa frase com uma afirmação de ``maior ciclagem de nutrientes'': esta última exige medidas adicionais. Explique também por que o valor-p do contraste único pode diferir do valor-p ajustado de Tukey para o mesmo par, embora a diferença estimada seja idêntica. Os procedimentos controlam famílias distintas de hipóteses; não escolha entre eles pelo menor valor-p depois de observar os resultados.
\clearpage
\section{Roteiro completo em R}
\scriptunidade{05_comparacoes.R}


\chapter{Regressão e gradientes ambientais}
\objetivo{Ajustar e interpretar regressão simples e múltipla; examinar resíduos e influência; distinguir incerteza da média, predição individual e associação causal.}

\section{Uma resposta ao longo de um gradiente}
A correlação resume a intensidade e a direção de uma associação, sem atribuir papéis de resposta e explicação. A regressão explicita esses papéis e modela a resposta média condicionada aos preditores. Na base de aves, a pergunta é: como a diversidade de Shannon se associa à heterogeneidade topográfica entre localidades?

A resposta $Y$ é Shannon, calculado com logaritmo natural. O preditor $X$ é o desvio-padrão da elevação, em metros, em uma janela espacial definida pelo protocolo. A comparação exige a mesma escala espacial e esforço amostral comparável. Os 36 locais são independentes no cenário simulado; em campo, essa condição precisa ser investigada a partir do desenho e dos processos ecológicos.

\begin{equation}Y_i=\beta_0+\beta_1X_i+\varepsilon_i.\end{equation}
$\beta_0$ é o valor médio previsto quando $X=0$; $\beta_1$ é a variação média da resposta para uma unidade adicional de $X$; $\varepsilon_i$ reúne desvios em torno da média condicional. A equação descreve uma relação estatística, não uma lei funcional determinística. Se $X=0$ estiver fora do domínio observado, o intercepto é necessário ao ajuste, mas pode ter interpretação ecológica limitada.

\section{Mínimos quadrados e qualidade do ajuste}
O método dos mínimos quadrados escolhe os coeficientes que minimizam $\sum_i(y_i-\widehat y_i)^2$. No modelo simples com intercepto:
\begin{equation}\widehat\beta_1=\frac{\sum_i(x_i-\bar x)(y_i-\bar y)}
{\sum_i(x_i-\bar x)^2},\qquad
\widehat\beta_0=\bar y-\widehat\beta_1\bar x.\end{equation}
O valor ajustado é $\widehat y_i=\widehat\beta_0+\widehat\beta_1x_i$ e o resíduo é $e_i=y_i-\widehat y_i$. Resíduo observado e erro populacional são conceitos relacionados, mas não idênticos: o primeiro depende dos coeficientes estimados na própria amostra.

\begin{equation}R^2=1-\frac{\sum_i e_i^2}{\sum_i(y_i-\bar y)^2},
\qquad RMSE=\sqrt{\frac{1}{n}\sum_i e_i^2}.\end{equation}
Com intercepto, $R^2$ descreve a fração da variação amostral em torno da média explicada pelo modelo ajustado. Não é a probabilidade de o modelo estar correto nem uma medida de causalidade. O RMSE está na escala da resposta. Aqui usamos divisor $n$ para o erro de ajuste; o erro-padrão residual do R usa os graus de liberdade, $n-p$, com $p$ coeficientes estimados, incluindo o intercepto.

Um modelo pode ajustar bem os dados usados na estimação e prever mal novos locais. O RMSE desta unidade é de ajuste, não de validação externa. Modelos destinados à previsão científica precisam de avaliação preditiva compatível com a estrutura espacial e temporal, além da inspeção de resíduos.

No ajuste simples executado, $\widehat H'=1{,}77292+0{,}02601T$, com $R^2=0{,}3716$ e RMSE de 0,2365. A inclinação corresponde a aproximadamente 0,026 unidade de Shannon por metro adicional de heterogeneidade, antes do ajuste pelos demais preditores.
\figura{u6_regressao.pdf}{.92}{Diversidade e topografia em 36 localidades simuladas. A faixa representa IC95\% pontual da resposta média; não é intervalo para uma nova observação.}

\section{Diagnósticos e limites}
A \textbf{linearidade} se refere à forma da média condicional; a \textbf{homocedasticidade}, à variância aproximadamente constante dos erros. A normalidade é relevante para a inferência clássica exata em pequenas amostras, não uma exigência de que cada preditor tenha distribuição normal. A independência não pode ser garantida por um Q--Q satisfatório.

Resíduos versus ajustados podem mostrar curvatura, grupos omitidos ou dispersão desigual. O Q--Q ajuda a avaliar caudas e assimetria. A distância de Cook combina discrepância e influência sobre o ajuste; um local extremo no preditor pode ter alta alavancagem mesmo com resíduo pequeno. Investigue erros, contexto e sensibilidade antes de excluir qualquer observação. A exploração estruturada ajuda a reconhecer problemas de variância, dependência e colinearidade antes da inferência \parencite{zuur2010}.

No exemplo, diferenças entre ambientes também contribuem para a resposta. A regressão simples é útil para introduzir a técnica, mas não isola a topografia de outros fatores. O modelo múltiplo a seguir incorpora ambiente e precipitação definidos pela pergunta; não seleciona automaticamente variáveis por valor-p.

\section{Confiança e predição}
O IC da média em $x_0$ descreve incerteza sobre $E(Y\mid X=x_0)$. O intervalo de predição de um novo local acrescenta a variação residual individual. Para o modelo simples normal:
\begin{align}
IC:\quad &\widehat y_0\pm t\,s_e\sqrt{h_0},\\
IP:\quad &\widehat y_0\pm t\,s_e\sqrt{1+h_0},
\end{align}
em que $s_e$ é o erro-padrão residual e $h_0=1/n+(x_0-\bar x)^2/\sum_i(x_i-\bar x)^2$. O quantil t usa o nível de confiança e os graus de liberdade correspondentes. Os intervalos são pontuais e condicionados à adequação do modelo.

Para heterogeneidade topográfica de 15 m, o modelo simples forneceu $\widehat H'=2{,}163$, IC95\% de 2,080 a 2,246 e IP95\% de 1,662 a 2,665. A predição é mais ampla porque inclui a variação entre novos locais. Prever fora do intervalo observado é extrapolar; mesmo uma reta bem ajustada pode falhar fora de sua faixa de suporte.

\section{Uma introdução à regressão múltipla}
Defina $P=(\text{precipitação}-2400)/100$: zero significa 2400 mm e uma unidade corresponde a 100 mm. Um modelo possível é:
\begin{equation}H'_i=\beta_0+\beta_1T_i+\beta_2P_i+
\gamma_1I_{M,i}+\gamma_2I_{R,i}+\varepsilon_i,\end{equation}
com $T$ representando topografia, $I_M$ um indicador de mosaico e $I_R$ de restauração. Conservado é a referência. O coeficiente de topografia passa a representar uma associação parcial, mantendo os demais preditores constantes no modelo. Combinações de condições com pouco suporte amostral precisam de cautela mesmo sem extrapolação univariada.

No ajuste executado, o coeficiente parcial foi 0,02687 por metro (IC95\% 0,01838 a 0,03535); o de chuva, 0,02165 por 100 mm. A comparação Mosaico menos Conservado foi $-0{,}4400$, ajustada por chuva e topografia. A inclusão desses termos não elimina todo confundimento nem converte o desenho observacional em experimento.

\begin{atencao}
Riqueza em função da precipitação e da sazonalidade é uma pergunta ecológica importante, mas riqueza é contagem. Um modelo linear generalizado pode ser mais apropriado que \texttt{lm()}, dependendo da distribuição, dispersão, esforço e desenho. Não copie o modelo de Shannon para riqueza sem revisar essas condições. Evitamos desenvolver modelos avançados nesta introdução.
\end{atencao}

\section{Exercícios da unidade}
\begin{enumerate}
\quest{1 | Em R}{Construa a dispersão entre Shannon e topografia. Descreva o padrão e a definição da variável explicativa.}
\quest{2 | Numérico e em R}{Obtenha equação ajustada, IC dos coeficientes, $R^2$ e RMSE. Interprete a inclinação e a limitação do intercepto.}
\quest{3 | Em R}{Examine os diagnósticos e identifique a observação de maior distância de Cook. Explique por que isso não autoriza exclusão automática.}
\quest{4 | Em R}{Para 15 m, compare IC da média e intervalo de predição. Verifique antes a faixa observada.}
\quest{5 | Em R e interpretação}{Ajuste o modelo múltiplo com chuva e ambiente. Interprete o coeficiente parcial de topografia e explique por que riqueza pode exigir outro modelo.}
\end{enumerate}
\subsection*{Oficina de interpretação do modelo}
Leia a tabela de coeficientes junto com os nomes das variáveis. No modelo múltiplo, chuva foi dividida por 100 depois de centrada; portanto, seu coeficiente não representa a variação por 1 mm. Para ambiente, cada coeficiente compara uma categoria com a referência Conservado, mantendo os demais preditores constantes.

Compare a inclinação simples com a parcial, seus intervalos e o diagnóstico de cada modelo. Uma mudança de coeficiente pode sinalizar diferenças de composição entre locais, mas não identifica sozinha um mecanismo causal. O roteiro também salva o diagnóstico do modelo múltiplo. Examine-o antes de usar a interpretação ajustada. Ao adaptar o código, revise o intervalo observado de cada preditor e não obtenha predições para combinações ambientais inexistentes na amostra apenas porque o R aceita o comando.
\clearpage
\section{Roteiro completo em R}
\scriptunidade{06_regressao.R}


\chapter[Estudo integrador: comunidades amazônicas]{Estudo integrador:\\comunidades amazônicas}
\objetivo{Percorrer uma análise reproduzível, distinguir descrição, comparação e ajuste; redigir conclusões ecológicas com estimativas, incerteza e limitações.}

\section{Pergunta, base e plano de análise}
O estudo didático pergunta como a diversidade de aves varia entre paisagens e se associa à topografia e à precipitação. A base reúne 36 localidades simuladas: 12 conservadas, 12 em mosaicos de uso do solo e 12 em restauração. Cada local contém o agregado de quatro pontos de escuta, com quatro horas totais em uma campanha. Não se utilizam os indivíduos ou pontos como réplicas de ambiente.

\begin{table}[H]\centering\small
\caption{Dicionário resumido da base de comunidades. Todas as linhas representam localidades e o grupo biológico é aves.}
\begin{tabular}{p{.39\linewidth}p{.51\linewidth}}\toprule
Coluna & Definição e unidade\\\midrule
\texttt{local}, \texttt{ambiente} & Identificação única e categoria de paisagem\\
\texttt{esforco\_horas} & Quatro horas por local\\
\texttt{riqueza}, \texttt{abundancia} & Contagens agregadas no local\\
\texttt{diversidade\_shannon} & Índice sintético com log natural\\
\texttt{precipitacao\_mm} & Precipitação anual, mm\\
\texttt{topografia\_m} & DP da elevação em janela comum, m\\
\texttt{sazonalidade\_cv} & CV da precipitação mensal, \%\\
\texttt{cobertura\_nativa\_pct} & Cobertura nativa no entorno, \%\\\bottomrule
\end{tabular}\end{table}

As variáveis aquáticas permanecem em bases separadas: não se atribui turbidez de igarapé a uma localidade de aves sem correspondência espacial e protocolo que justifiquem essa integração. Da mesma forma, uma matriz por espécie teria estrutura diferente da tabela com índices por local. Não há uma coluna de espécie na base agregada porque cada linha reúne a comunidade.

O plano é: auditar identificadores e esforço; descrever Shannon por ambiente; visualizar a distribuição; comparar médias não ajustadas; examinar uma regressão com topografia, chuva e ambiente; redigir os resultados. A sequência não deve ser usada para procurar repetidamente um resultado significativo. As perguntas e a família de comparações precisam estar definidas.

\section{Explorar sem antecipar causalidade}
Médias e desvios-padrão resumem os grupos; quartis mostram a distribuição central. Os pontos no boxplot permitem identificar sobreposição e heterogeneidade. Não interprete um único local de diversidade alta como evidência de que todo o ambiente favorece diversidade. Também não descreva a categoria de paisagem como tratamento casualizado.

\figura{u7_ambientes.pdf}{.90}{Distribuição de Shannon por ambiente em localidades simuladas. Os contrastes descrevem associações entre ambientes e resposta.}

O esforço é igual, mas a detectabilidade pode variar entre ambientes. Em dados reais, disponibilidade de habitat, observador, época e acessibilidade podem influenciar tanto o grupo amostrado quanto os resultados. A amostragem probabilística, protocolos padronizados e modelos apropriados à observação precisam ser considerados antes da inferência regional.

\section{Comparar médias e ajustar associações}
A ANOVA de Shannon por ambiente testa a igualdade de médias \textbf{não ajustadas} por clima ou topografia. A estatística F pode ser usada em um estudo observacional sob as suposições do modelo; o que muda é a interpretação causal. O termo ``entre grupos'' é mais adequado que atribuir um efeito experimental ao ambiente.

Após avaliar os diagnósticos, Tukey descreve as diferenças da família de três comparações. A regressão múltipla adiciona precipitação e topografia e estima contrastes condicionais aos preditores incluídos. Ela não responde exatamente à mesma pergunta que a ANOVA simples: uma comparação é marginal aos demais fatores na amostra; a outra é ajustada pelo modelo.

\begin{exemplo}
No modelo múltiplo, a equação estimada, com chuva centrada e escalonada, foi aproximadamente:
\[
\widehat H'=1{,}9746+0{,}02687T+0{,}02165P
-0{,}4400I_M-0{,}2884I_R.
\]
Para Conservado, $I_M=I_R=0$; para Mosaico, $I_M=1$ e $I_R=0$; para Restauração, o inverso. Essas diferenças são estimativas para dados simulados, condicionadas à estrutura escolhida.
\end{exemplo}

Cobertura e sazonalidade estão disponíveis para ampliar a discussão, mas não são incluídas automaticamente. Mais preditores consomem informação e podem representar variáveis correlacionadas, mecanismos intermediários ou escalas incompatíveis. A escolha exige hipótese, suporte amostral e avaliação da colinearidade. Não se recomenda testar todas as combinações e apresentar apenas a melhor por valor-p.

\section{Como interpretar resultados estatísticos em estudos ecológicos?}
Comece pela unidade de análise e pelo tamanho amostral efetivo. Identifique o desenho e o método. Apresente a estimativa na escala original, com intervalo e, quando pertinente, estatística de teste, graus de liberdade e valor-p. Relacione a magnitude à pergunta ecológica e explicite as limitações que podem alterar sua interpretação.

Uma redação possível é: ``Em 36 localidades simuladas, Shannon apresentou associação parcial positiva com heterogeneidade topográfica de 0,0269 por metro (IC95\% 0,0184 a 0,0354), após ajuste por chuva e ambiente. Como os ambientes foram observados, o resultado não estabelece causalidade.'' O script produz essa frase a partir dos coeficientes realmente ajustados.

Evite frases como ``o valor-p prova o efeito'', ``não houve diferença porque $p>0{,}05$'' ou ``o ambiente explicou 100\% do processo''. Ausência de evidência não demonstra equivalência; um efeito estatisticamente detectável pode ser pequeno para a decisão; uma associação pode refletir variáveis omitidas. As recomendações de comunicação da ASA enfatizam que conclusões não devem depender somente de um limiar \parencite{wasserstein2016}.

\begin{dica}
Para transferir o roteiro à dissertação, escreva um parágrafo de métodos antes de executar: quais locais representam a população, como foram selecionados, qual é a unidade, quais medições se repetem, quais preditores respondem à hipótese e qual estimativa sustenta a conclusão. Se esse parágrafo estiver indefinido, o comando estatístico ainda não está pronto para uso científico.
\end{dica}

\section{Exercícios da unidade}
\begin{enumerate}
\quest{1 | Conceitual e em R}{Audite a base, confirme 12 locais por ambiente e esforço constante. Explique a unidade e uma limitação de detectabilidade.}
\quest{2 | Em R}{Calcule médias, DP e quartis de Shannon por ambiente e produza o gráfico com observações individuais.}
\quest{3 | Em R e interpretação}{Ajuste ANOVA e Tukey, examine diagnósticos e interprete um contraste como associação observacional.}
\quest{4 | Em R}{Ajuste Shannon por topografia, chuva e ambiente. Compare o significado dos coeficientes com o da ANOVA não ajustada.}
\quest{5 | Síntese}{Redija um resultado com tamanho amostral, coeficiente parcial, IC e limite causal. Indique uma modificação necessária para aplicar o fluxo a dados repetidos ou espacialmente agrupados.}
\end{enumerate}
\subsection*{Produto final da atividade}
Entregue um relatório curto com a pergunta, o desenho e a unidade de análise; uma tabela de médias e dispersão; um gráfico com observações individuais; um resultado de comparação; e a interpretação do coeficiente parcial de topografia. Acrescente uma limitação que derive do desenho, e não apenas de um aviso genérico de software.

Anexe o script executado e o registro da sessão. Explique por que a comparação não ajustada entre ambientes e a regressão múltipla respondem a perguntas diferentes. A conclusão deve permanecer proporcional à evidência: dados simulados demonstram um procedimento; dados observacionais reais podem revelar associações; uma afirmação causal exige sustentação adicional do desenho e das hipóteses. A frase gerada pelo script é um ponto de partida a revisar, não uma conclusão automática para qualquer base.
\clearpage
\section{Roteiro completo em R}
\scriptunidade{07_integrador.R}


\backmatter
\fancyhead[R]{\small\textcolor{UEAPDarkGreen}{Material de consulta}}
\chapter{Gabarito comentado}
As respostas numéricas abaixo correspondem aos dados e à semente distribuídos nesta edição. Ao alterar a base, reexecute os scripts e não conserve estes resultados como se fossem universais. As resoluções completas em R aparecem ao final das respectivas unidades.

\section*{Unidade 1}
\begin{enumerate}
\item Ambiente: qualitativa nominal; riqueza: discreta; precipitação: contínua; local: identificador. Parâmetro descreve a população; estimativa é o valor amostral de um estimador.
\item São 36 localidades na análise entre locais, não 144 réplicas independentes. O agregado de pontos define a resposta de cada local.
\item A base não contém ausências ou identificadores duplicados e respeita $H'\leq\ln S$. Coerência interna não comprova seleção representativa nem detectabilidade constante.
\item A tabela exportada tem uma linha por local, quatro horas de esforço e 12 locais por ambiente. A inferência regional dependeria de como os locais reais fossem selecionados.
\end{enumerate}

\section*{Unidade 2}
\begin{enumerate}
\item Média 4, mediana 2 e moda 2 mg/L. Média ponderada do segundo exemplo: 2,5 mg/L. A média é mais sensível ao valor 12.
\item Amplitude 11 mg/L; variância 20,5 (mg/L)$^2$; DP 4,528 mg/L; EP 2,025 mg/L; CV 113,19\%; quartis 2, 2 e 3 mg/L; AIQ 1 mg/L; P90 8,4 mg/L pelo tipo 7.
\item A média de oxigênio é 6,448 mg/L em Conservado e 5,465 mg/L em Mosaico. Os decis de turbidez são posições empíricas dos 24 igarapés, não limites ambientais normativos; consulte os valores completos no log da unidade 2.
\item Histograma e densidade evidenciam a distribuição; boxplot compara grupos; dispersão apresenta pares. Escolhas gráficas alteram a aparência e não estabelecem causalidade nem independência.
\end{enumerate}

\section*{Unidade 3}
\begin{enumerate}
\item $H_0:\mu_C-\mu_M=0$ e $H_1:\mu_C-\mu_M\ne0$. O valor-p é condicional à nula e ao modelo, não uma probabilidade posterior da hipótese.
\item Diferença 0,983 mg/L, IC95\% 0,578 a 1,389 mg/L; Welch $t=5{,}047$, $gl\approx20{,}468$, $p\approx0{,}000058$. A diferença não constitui efeito causal do ambiente.
\item Q--Q e Shapiro avaliam aspectos da distribuição, com limitações de tamanho amostral; não comprovam independência. Um $p>0{,}05$ pode refletir baixa sensibilidade do teste.
\item O teste pareado usa depois menos antes, com $n=12$ diferenças e 11 gl. Média 0,2975 mg/L, IC95\% 0,1360 a 0,4590; $t=4{,}055$, $p=0{,}0019$. O controle do pareamento não elimina confundimento temporal.
\item O planejamento retorna $n=12{,}570$ por grupo; arredondando para cima, 13 por grupo, total 26. O cálculo pressupõe grupos independentes e uma variância comum definida no planejamento, sem perdas ou agrupamento.
\end{enumerate}

\section*{Unidade 4}
\begin{enumerate}
\item Tratamento: umidade; resposta: perda de massa em \%; unidade: recipiente independente. São oito unidades por nível no DIC.
\item $SQ_T\approx1631{,}6$, $SQ_{\mathrm{trat}}\approx1248{,}1$, $SQ_E\approx383{,}4$; $F\approx34{,}18$. Os gl são $a-1=2$, $N-a=21$ e $N-1=23$.
\item Padrões residuais podem sugerir forma inadequada, variância desigual ou influência. Fligner é auxiliar para variâncias; a independência precisa do protocolo e da estrutura de dependência.
\item DBC: 18 unidades, seis blocos completos, gl 5 para bloco, 2 para tratamento e 10 para resíduo. No ajuste, $F_{\mathrm{trat}}\approx105{,}43$. Uma observação por célula não permite separar interação e erro puro.
\end{enumerate}

\section*{Unidade 5}
\begin{enumerate}
\item F global rejeita igualdade de todas as médias; não identifica todos os pares distintos. Tukey controla a família de comparações par a par.
\item Alta menos Baixa: 17,624 pontos percentuais, IC simultâneo 12,239 a 23,009; $p$ ajustado aproximadamente $1{,}47\times10^{-7}$. Diferença entre porcentagens não equivale a aumento relativo percentual.
\item IC de médias individuais representam outra quantidade. Não infira o teste de diferença apenas por sua sobreposição visual.
\item O contraste $(-1,0,1)$ tem estimativa 17,624 e EP aproximadamente 2,136 pontos percentuais. A justificativa de contraste único exige formulação prévia; uma família exige estratégia de multiplicidade.
\end{enumerate}

\section*{Unidade 6}
\begin{enumerate}
\item O gráfico sugere associação positiva. A topografia é DP da elevação em metros, não altitude média. A associação simples pode refletir outros preditores.
\item Equação simples: $\widehat H'=1{,}77292+0{,}02601T$; IC da inclinação 0,01422 a 0,03780; $R^2=0{,}3716$; RMSE 0,2365. A inclinação é a variação média de Shannon por metro. O intercepto em zero está fora do domínio observado (4,4 a 30 m).
\item A identificação por \texttt{which.max(cooks.distance())} retorna a linha 4, local L04. Não constitui critério automático de exclusão. Verifique dados, contexto e sensibilidade.
\item Para 15 m: média ajustada 2,163; IC 2,080 a 2,246; IP 1,662 a 2,665. O IP inclui variabilidade de uma nova observação.
\item Coeficiente parcial de topografia: 0,02687 (IC 0,01838 a 0,03535), ajustado por chuva e ambiente. Riqueza exige considerar contagens, distribuição e esforço; não há transferência automática do modelo normal.
\end{enumerate}

\section*{Unidade 7}
\begin{enumerate}
\item São 36 locais únicos, com esforço igual. Detecção ainda pode diferir entre ambientes em um estudo real.
\item Médias (DP): Conservado 2,417 (0,272); Mosaico 1,960 (0,242); Restauração 2,178 (0,211). Quartis: Conservado 2,246/2,385/2,500; Mosaico 1,800/1,913/2,020; Restauração 2,064/2,121/2,345. Pontos individuais mostram variação e sobreposição.
\item ANOVA: $F_{2,33}=10{,}65$, $p=0{,}00027$. Mosaico menos Conservado: $-0{,}4575$ (IC simultâneo $-0{,}7009$ a $-0{,}2141$). Os outros dois IC incluem zero. A comparação é observacional e não ajustada por clima ou topografia; não demonstra efeito causal.
\item O modelo múltiplo estima diferenças entre ambientes condicionadas à chuva e à topografia. Mosaico menos Conservado: $-0{,}4400$; Restauração menos Conservado: $-0{,}2884$. Os contrastes não são efeitos causais.
\item A frase gerada em \texttt{resultados/u7\_redacao.txt} inclui 36 locais, coeficiente e IC. Para dados repetidos ou agrupados, reveja o tamanho amostral efetivo e utilize uma estrutura que modele essa dependência; adicionar linhas ao mesmo modelo não resolve o problema.
\end{enumerate}

\printbibliography[heading=bibintoc,title={Referências e leituras de apoio}]
\end{document}
